% Figure: plasma sheath. Potential drops from the quasineutral bulk to a
% negative wall over a few Debye lengths; the schematic shows the bulk,
% the sheath edge, and the ion-acceleration region that produces the
% directional bombardment used in reactive-ion etching.
\begin{figure}[htbp]
\centering
\begin{tikzpicture}
\begin{axis}[
  width=0.78\textwidth, height=0.46\textwidth,
  xlabel={distance from wall (Debye lengths $x/\lambda_D$)},
  ylabel={potential $\phi$ (V)},
  domain=0:8, samples=120,
  xmin=0, xmax=8, ymin=-22, ymax=3,
  grid=both, grid style={enggray!20},
  every axis plot/.append style={thick},
]
  % sheath: potential rises from wall (-20 V) to ~0 in the bulk
  \addplot[engblue] {-20*exp(-x)};
  \draw[dashed, enggray] (axis cs:0,0) -- (axis cs:8,0);
  \node[engblue, font=\scriptsize, anchor=west] at (axis cs:3.2,-1.6)
    {quasineutral bulk $\phi\approx 0$};
  \draw[engred, ->, thick] (axis cs:1.6,-15) -- (axis cs:0.2,-19)
    node[midway, above, sloped, font=\scriptsize] {ion acceleration};
  \node[font=\scriptsize, anchor=south] at (axis cs:0.0,-21.5) {wall};
  \node[font=\scriptsize, engred, anchor=south west] at (axis cs:0.05,-21.5)
    {$\;V_{\mathrm{wall}}<0$};
\end{axis}
\end{tikzpicture}
\caption[Plasma sheath potential]{The plasma sheath. Because electrons are
far more mobile than ions, any surface in contact with a plasma charges
negative until it repels enough electrons to balance the ion and electron
fluxes. The resulting potential drop is confined to a sheath a few Debye
lengths thick at the wall, while the bulk plasma stays quasineutral at
nearly constant potential. Ions entering the sheath are accelerated
toward the wall along the field gradient, producing the directional
bombardment that gives reactive-ion etching its anisotropy. The
exponential profile shown is a schematic; the exact sheath solution
follows the Child-Langmuir law in the high-voltage limit.}
\label{fig:vol04ch13:sheath-potential}
\end{figure}
